\documentclass[12pt,a4paper]{article}
\usepackage[utf8]{inputenc}
\usepackage[slovak]{babel}
\usepackage{graphicx}
\usepackage{pgf-pie}
\usepackage{geometry}
\usepackage{amsmath}
\usepackage{hyperref}
\geometry{margin=2.5cm}
\usepackage{pgfplots}
\pgfplotsset{compat=1.18}
\usepackage{float}

\title{Analýza odporúčacieho systému pre eventy}
\author{Marma}
\date{\today}

\begin{document}

\section{Testovanie výsledkov}

\subsection{Metodológia a testovacie dáta}
V rámci testovania sme simulovali prostredie so 100 používateľmi a 200 udalosťami. Na overenie správnosti algoritmu sme zvolili dva prístupy:
\begin{enumerate}
    \item \textbf{Analýza vplyvu váh:} Sledovanie zmien odporúčaní pre jedného používateľa pri zmene preferencií a overenie správnosti výsledkov.
    \item \textbf{Analýza personalizácie:} Porovnanie odporúčaní pre dvoch používateľov s odlišnou históriou pri identických váhach.
\end{enumerate}


\section{Citlivosť na zmenu váh (Používateľ Janko)}
\subsection{Navštívené eventy}

Používateľ Janko navštívil celkovo 8 udalostí. Ich rozdelenie podľa hlavných kategórií je spracované v tabuľke \ref{tab:rozdelenie_kat}.

\begin{table}[H]
\centering
\caption{Rozdelenie návštev podľa kategórií}
\label{tab:rozdelenie_kat}
\begin{tabular}{|l|c|}
\hline
\textbf{Kategória} & \textbf{Počet návštev} \\ \hline
Sports & 4 \\
Education & 3 \\
Business \& Tech & 1 \\ \hline
\textbf{Spolu} & \textbf{8} \\ \hline
\end{tabular}
\end{table}

Podrobný zoznam konkrétnych udalostí, ktoré Janko v minulosti navštívil, uvádzame v tabuľke \ref{tab:konkretne_eventy}.

\begin{table}[H]
\centering
\caption{Konkrétne navštívené udalosti a ich kategorizácia}
\label{tab:konkretne_eventy}
\small % Zmenšili sme písmo, aby tabuľka nebola príliš široká
\begin{tabular}{|l|l|l|}
\hline
\textbf{Názov udalosti} & \textbf{Podkategória} & \textbf{Kategória} \\ \hline
udalost 15 \#89  & Pitch Night              & Business \& Tech \\
udalost 9 \#173  & Workshop (ručné práce)   & Education \\
udalost 4 \#123  & Osobný rozvoj            & Education \\
udalost 7 \#171  & Strength \& Conditioning & Sports \\
udalost 9 \#128  & Team sports              & Sports \\
udalost 3 \#92   & Winter sports            & Sports \\
udalost 10 \#39  & Osobný rozvoj            & Education \\
udalost 12 \#146 & Recreational sports      & Sports \\ \hline
\end{tabular}
\end{table}

\subsection{Obľúbené eventy}

Okrem histórie navštívených podujatí systém berie do úvahy aj udalosti, ktoré používateľ explicitne označil ako obľúbené. Tieto dáta majú vo výpočtovom vzorci vyššiu váhu vďaka faktoru $B_{fav}$ (Obľúbený bonus). V tabuľke \ref{tab:oblubene_udalosti} uvádzame zoznam týchto podujatí pre používateľa Janka.

\begin{table}[H]
\centering
\caption{Udalosti označené používateľom ako obľúbené}
\label{tab:oblubene_udalosti}
\begin{tabular}{|l|l|l|}
\hline
\textbf{Názov udalosti} & \textbf{Podkategória} & \textbf{Kategória} \\ \hline
udalost 7 \#171  & Strength \& Conditioning & Sports \\ \hline
\end{tabular}
\end{table}



\subsection{Vizuálna analýza histórie používateľa Janka}

Na obrázku \ref{fig:porovnanie_grafov} môžeme vidieť hĺbkovú analýzu Jankovho profilu. Zatiaľ čo ľavý graf zobrazuje dominanciu športu na úrovni hlavných kategórií, pravý graf odhaľuje, že jeho záujmy sú v rámci podkategórií veľmi diverzifikované.

\begin{figure}[H]
\centering

% --- PRVÝ GRAF (Hlavné kategórie) ---
\begin{minipage}{0.48\textwidth}
    \centering
    \begin{tikzpicture}[scale=0.7]
        % Nastavenie písma priamo pre uzly v tikz
        \tikzset{every node/.style={font=\scriptsize}}
        \pie[text=legend, 
             radius=1.8, 
             color={blue!60, green!60, orange!60}]{
            50/Sports,
            37.5/Education,
            12.5/Business \& Tech
        }
    \end{tikzpicture}
    \caption{Hlavné kategórie}
\end{minipage}
\hfill
% --- DRUHÝ GRAF (Podkategórie) ---
\begin{minipage}{0.48\textwidth}
    \centering
    \begin{tikzpicture}[scale=0.7]
        \tikzset{every node/.style={font=\scriptsize}}
        \pie[text=legend, 
             radius=1.8, 
             color={cyan!60, yellow!60, red!60, gray!60, violet!60, brown!60, green!60}]{
            25/Osobný rozvoj,
            12.5/Pitch Night,
            12.5/Workshop,
            12.5/Strength \& Cond.,
            12.5/Team sports,
            12.5/Winter sports,
            12.5/Recreational sports
        }
    \end{tikzpicture}
    \caption{Podkategórie}
\end{minipage}

\caption{Porovnanie preferencií používateľa Janka na úrovni kategórií a podkategórií}
\label{fig:porovnanie_grafov}
\end{figure}




\section{Testovanie odporúčacieho systému}



Pre výpočet finálneho skóre jednotlivých podujatí používame vážený súčet faktorov definovaný v kapitole [Este doplnitcislo] (vzťah \ref{eq:score_formula}):

\begin{equation}
Score = w_1 \cdot C_{main} + w_2 \cdot C_{sub} + w_3 \cdot D + w_4 \cdot R + w_5 \cdot P + w_6 \cdot B_{fav}
\label{eq:score_formula}
\end{equation}

Tento prístup umožňuje systému adaptovať sa na rôzne typy používateľov prostredníctvom úpravy váhového vektora. Pre potreby validácie sme definovali sedem testovacích scenárov, v ktorých systematicky meníme váhy $w_1$ až $w_6$. Týmto spôsobom demonštrujeme citlivosť systému na špecifické preferencie (napr. uprednostnenie geografickej blízkosti pred historickou relevanciou).

Každý test obsahuje konfiguráciu váh, prehľad TOP 5 odporúčaných udalostí a vizuálny rozbor príspevkov jednotlivých faktorov k výslednému skóre.


testovania 

\subsubsection{Analýza vplyvu váh pri rovnomernom rozdelení (Test 1)}

Pri teoreticky rovnomernom rozdelení váh (každý faktor $w_i \approx 16.67\,\%$) sa v popredí umiestňujú najmä udalosti z obľúbených kategórií s vysokým ratingom. Tieto konkrétne výsledky a presné číselné hodnoty pre všetkých 200 testovaných udalostí v rámci simulácie pozorujeme v \textbf{Prílohe XY1}, kde je spracovaný kompletný výstup algoritmu pre tento scenár.

Tento jav je spôsobený vlastnosťami normalizovaných vstupov, ktoré pri rovnomernom váhovaní ovplyvňujú finálne skóre nasledovne:

\begin{itemize}
    \item \textbf{Obľúbený bonus ($w_6$):} Nadobúda binárne hodnoty (0 alebo 1). Pri rovnakých váhach tak automaticky posúva obľúbených organizátorov o celú váhovú kategóriu vyššie oproti ostatným.
    \item \textbf{Rating ($w_4$):} Keďže väčšina podujatí v testovacej sade má vysoké hodnotenie (blízke 5 hviezdičkám), normalizovaná hodnota $R$ sa stabilne drží blízko 1.0, čím vytvára silný a konštantný tlak na skóre.
    \item \textbf{Popularita ($w_5$):} Výrazne rastie až pri napĺňaní kapacity. Je limitovaná hodnotou 1.0, ktorú však v praxi nedosiahne, nakoľko plne obsadené podujatia systém z ponuky vyraďuje.
    \item \textbf{Vzdialenosť ($w_3$):} Prirodzene nedosahuje hodnotu 1.0, pretože vyžaduje nulovú vzdialenosť od používateľa. Jej vplyv je stredne silný a závisí od hustoty podujatí v okolí.
    \item \textbf{Kategorizácia ($w_1, w_2$):} Vykazujú vysokú mieru zotrvačnosti. Kvôli diverzifikovanej histórii používateľa (napr. profil v Prílohe B so záujmami rozdelenými medzi Business, Art a Music) sú ich normalizované hodnoty nízke. Pri nízkej alebo rovnomernej váhe sú tieto rozdiely „prekryté“ faktormi s prirodzene vyššími hodnotami (rating, vzdialenosť).
\end{itemize}



Z uvedeného vyplýva dôležitý poznatok: **čím viac chceme personalizovať výsledky podľa histórie používateľa, tým agresívnejšie musíme zvyšovať váhy kategórií.** Pri nízkych váhach sa tematické rozdiely prejavujú minimálne. Pre dosiahnutie "subjektívnej rovnováhy" je preto nevyhnutné upraviť váhy tak, aby faktory s nízkym rozptylom (rating, obľúbenosť) mali nižšiu váhu a faktory s vysokou variabilitou a nižšími vstupnými hodnotami (kategórie) váhu výrazne vyššiu.
% ───────────────────────────────────────────────────────────────────────────────
\subsubsection{Porovnanie vplyvu faktorov pri extrémnej váhe}
% ───────────────────────────────────────────────────────────────────────────────

Tabuľka \ref{tab:vplyv_faktorov} sumarizuje reálny dopad jednotlivých faktorov na výsledné skóre v prípadoch, kedy im bola pridelená dominantná váha 60 \%.

\begin{table}[H]
\centering
\small
\begin{tabular}{|l|c|p{9cm}|}
\hline
\textbf{Faktor} & \textbf{Reálny vplyv*} & \textbf{Zdôvodnenie správania} \\
\hline
Obľúbený bonus & 76 \% skóre & Binárny charakter (0/1) spôsobuje najprudšie zmeny v poradí. \\
\hline
Rating & 74 \% skóre & Vysoký priemer hodnôt $R \approx 0.95$ spôsobuje takmer fixný prírastok k skóre. \\
\hline
Vzdialenosť & 64 \% skóre & Vysoký rozptyl medzi blízkymi (0.85) a vzdialenými (0.35) udalosťami. \\
\hline
Kategórie & 30--41 \% skóre & Závisí od histórie; pri pestrej histórii sú hodnoty nízke. \\
\hline
Popularita & 31 \% skóre & V testovacej sade najslabší faktor kvôli nízkej obsadenosti podujatí. \\
\hline
\end{tabular}
\caption{Vplyv faktorov pri 60\% konfigurácii (*priemerný podiel na skóre TOP 5 udalostí)}
\label{tab:vplyv_faktorov}
\end{table}



% ───────────────────────────────────────────────────────────────────────────────
\subsubsection{Kritické zhodnotenie faktora Podkategórie}
% ───────────────────────────────────────────────────────────────────────────────

Počas testovania sa ukázalo, že \textbf{podkategória} má najnižšiu tendenciu ovplyvňovať výsledok. Jej zachovanie v algoritme má však strategický význam:
\begin{itemize}
    \item \textbf{Implicitné učenie:} Systém reaguje na zmenu záujmov skôr, než si používateľ manuálne upraví profil.
    \item \textbf{Jemnosť odporúčaní:} Zabraňuje tomu, aby bol používateľ "zaplavený" všetkými športovými udalosťami, ak v skutočnosti preferuje len jogu.
\end{itemize}

% ───────────────────────────────────────────────────────────────────────────────
\subsubsection{Odporúčané produkčné nastavenie}
% ───────────────────────────────────────────────────────────────────────────────

Na základe analýzy navrhujeme pre reálne nasadenie nasledovné rozdelenie váh, ktoré uprednostňuje tematickú zhodu pred čistou popularitou či geografickou blízkosťou.

\begin{table}[H]
\centering
\begin{tabular}{|l|c|l|}
\hline
\textbf{Faktor} & \textbf{Váha} & \textbf{Zdôvodnenie} \\
\hline
Hlavná kategória ($w_1$) & 30\% & Prioritná tematická zhoda. \\
Podkategória ($w_2$) & 25\% & Jemné doladenie podľa špecifických záujmov. \\
Vzdialenosť ($w_3$) & 15\% & Zabezpečenie praktickej dostupnosti. \\
Rating ($w_4$) & 10\% & Filter kvality. \\
Obľúbený bonus ($w_6$) & 10\% & Priama personalizácia. \\
Popularita ($w_5$) & 10\% & Podpora trendových podujatí. \\
\hline
\textbf{Súčet} & \textbf{100\%} & \\
\hline
\end{tabular}
\caption{Navrhované produkčné nastavenie váhového vektora}
\label{tab:produkcia}
\end{table}


\newpage
\appendix
\section{Detailné výsledky testovacích scenárov}
\input{GrafyRecomandationSystemu.tex}




\end{document}
